\subsection{滤子的基}

若 \(\mathfrak{G}\) 是 X 上滤子 \(\mathfrak{F}\) 的一个子基（第 2 目），则 \(\mathfrak{F}\) 一般并不是 X 中包含 \(\mathfrak{G}\) 中某个集合的子集所成之集；为使 \(\mathfrak{G}\) 具有这一性质，必须且只需 \(\mathfrak{G}\) 中集合的每个有限交都包含 \(\mathfrak{G}\) 的某个集合。于是有下列命题：

命题 2. 设 \(\mathfrak{B}\) 是集合 X 的子集所成之集。则 X 中包含 \(\mathfrak{B}\) 中某个集合的子集所成之集是一个滤子，当且仅当 \(\mathfrak{B}\) 具有下列两条性质：

（B \(_{I}\)）B 中两个集合的交包含 B 的某个集合。

（B \(_{II}\)）B 非空，且 X 的空子集不属于 B。

定义 3. 集合 X 的满足公理（B \(_{I}\)）与（B \(_{II}\)）的子集所成之集 B 称为它所生成的滤子的基。若两个滤子基生成同一个滤子，则称它们等价。

若 \(\mathfrak{G}\) 是滤子 \(\mathfrak{F}\) 的一个子基，则集合 \(\mathfrak{G}'\)（\(\mathfrak{G}\) 中集合的有限交所成之集）是 \(\mathfrak{F}\) 的一个基（第 2 目）。

命题 3. 子集 \(\mathfrak{B}\)——滤子 \(\mathfrak{F}\)（\(\mathbf{X}\) 上的）的子集——是 \(\mathfrak{F}\) 的基，当且仅当 \(\mathfrak{F}\) 的每个集合都包含 \(\mathfrak{B}\) 的某个集合。

若 \(\mathfrak{B}\) 是 \(\mathfrak{F}\) 的基，则显然 \(\mathfrak{F}\) 的每个集合都包含 \(\mathfrak{B}\) 的某个集合；反之，若 \(\mathfrak{F}\) 的每个集合都包含 \(\mathfrak{B}\) 的某个集合，则 \(X\) 中包含 \(\mathfrak{B}\) 的某个集合的子集所成之集与 \(\mathfrak{F}\) 一致（由 \((\mathbf{F}_{\mathbf{I}})\)）。

例如，Fréchet 滤子（第 1 目）是被视为按关系 \(\leqslant\) 定向的有序集 N 的截面滤子。设 \(\mathfrak{F}\) 是集合 Z 上的一个滤子。由于 \(\mathfrak{F}\) 按关系 \(\supset\) 是有向的{[}由公理 \((\mathrm{F}_{\mathrm{II}})\){]}，我们可以在 \(\mathfrak{F}\) 上定义截面滤子；这里 \(\mathfrak{F}\) 相对于集合 A\(\in\mathfrak{F}\) 的截面是所有 M\(\in\mathfrak{F}\)（满足 M\(\subset\)A 者）所成之集 S(A)。这个滤子称为滤子 \(\mathfrak{F}\) 的截面滤子。

命题 4. 在集合 X 上，滤子 \(\mathfrak{F}'\)（以 \(\mathfrak{B}'\) 为基的）细于滤子 \(\mathfrak{F}\)（以 \(\mathfrak{B}\) 为基的），当且仅当 \(\mathfrak{B}\) 的每个集合都包含 \(\mathfrak{B}'\) 的某个集合。

这是定义 2 和定义 3 的直接推论。

推论. 集合 X 上两个滤子基 B、B' 等价，当且仅当 B 的每个集合都包含 B' 的某个集合，且 B' 的每个集合都包含 B 的某个集合。

滤子基的例子. 1) 设 X 是拓扑空间。命题 3 表明，一点 \(x \in X\) 的邻域滤子的基恰好就是 \(x\) 的基本邻域系（§ 1，第 3 目，定义 5）。

\begin{enumerate}
\def\labelenumi{\arabic{enumi})}
\setcounter{enumi}{1}
\tightlist
\item
  设 X 是关于关系 \((\sigma)\) 的非空有向集（《集合论》，第三章，§ 1，第 10 目）。对每个 \(a \in X\)，集合 \(S(a)\)——由所有这样的 \(x \in X\)（满足 \(a(\sigma)x\) 者）组成——称为 X 相对于元素 a 的截面。于是 X 的截面所成之集 S 是一个滤子基，因为它显然满足 \((\mathbf{B}_{\mathbf{II}})\)，且若 a、b 是 X 的任意两个元素，则按假设存在元素 \(c \in X\) 使得 \(a(\sigma)c\) 与 \(b(\sigma)c\)，从而
\end{enumerate}

\[
\mathbf {S} (c) \subset \mathbf {S} (a) \cap \mathbf {S} (b),
\]

于是 \((\mathbf{B}_{\mathbf{I}})\) 得到满足。由 S 生成的滤子称为有向集 X 的截面滤子。

